% !TEX root = ../main.tex
% Moved from Background (Vehicle model) on 2026-11-25. Background keeps one sentence on the
% notation (Simulation, BlueROV2 Heavy in MarineGym) and points here.
\section{Vehicle model}
\label{app:vehicle-model}
A common approach to modelling underwater vehicles is based on Fossen's 6-DOF equations of motion \parencite{fossen2021handbook}. The motion is described using two reference frames: an Earth-fixed North-East-Down (NED) frame $\{n\}$, which is treated as inertial, and a body-fixed frame $\{b\}$ with its origin at the vehicle centre and with $x_b$ pointing forward, $y_b$ to starboard and $z_b$
down. The state of the vehicle is given by
\begin{beqn}
\begin{aligned}
\boldsymbol{\eta} &= [x,\, y,\, z,\, \phi,\, \theta,\, \psi]^\top \in \mathbb{R}^6, \\
\boldsymbol{\nu}  &= [u,\, v,\, w,\, p,\, q,\, r]^\top \in \mathbb{R}^6,
\end{aligned}
\end{beqn}
where $\boldsymbol{\eta}$ is the position $(x, y, z)$ and orientation, given as the Euler angles roll $\phi$, pitch $\theta$ and yaw $\psi$, of the vehicle expressed in $\{n\}$. $\boldsymbol{\nu}$ is the vector of linear velocities (surge $u$, sway $v$, heave $w$) and angular velocities (roll rate $p$, pitch rate $q$, yaw rate $r$) expressed in $\{b\}$. The two are related by the kinematic equation
\begin{beqn}
\begin{aligned}
\dot{\boldsymbol{\eta}} = \mathbf{J}_\Theta(\boldsymbol{\eta})\,\boldsymbol{\nu},
\qquad
\mathbf{J}_\Theta(\boldsymbol{\eta}) =
\begin{bmatrix}
\mathbf{R}^n_b(\boldsymbol{\Theta}) & \mathbf{0}_{3\times 3} \\
\mathbf{0}_{3\times 3} & \mathbf{T}_\Theta(\boldsymbol{\Theta})
\end{bmatrix} \in \mathbb{R}^{6\times 6},
\end{aligned}
\end{beqn}
where $\boldsymbol{\Theta} = [\phi, \theta, \psi]^\top$, $\mathbf{R}^n_b \in \mathbb{R}^{3\times 3}$ is the rotation matrix from $\{b\}$ to $\{n\}$, and $\mathbf{T}_\Theta \in \mathbb{R}^{3\times 3}$ maps the body angular rates to Euler-angle rates.

The dynamics are given by
\begin{beqn}
\begin{aligned}
\mathbf{M}\dot{\boldsymbol{\nu}}
+ \mathbf{C}(\boldsymbol{\nu})\boldsymbol{\nu}
+ \mathbf{D}(\boldsymbol{\nu})\boldsymbol{\nu}
+ \mathbf{g}(\boldsymbol{\eta})the FAO primary source for the production figures;
= \boldsymbol{\tau},
\end{aligned}
\end{beqn}
where $\boldsymbol{\tau} \in \mathbb{R}^6$ is the vector of control forces and
moments produced by the thrusters, and
\begin{beqn}
\begin{aligned}
\mathbf{M} &= \mathbf{M}_{RB} + \mathbf{M}_A \in \mathbb{R}^{6\times 6}, \\
\mathbf{C}(\boldsymbol{\nu}) &= \mathbf{C}_{RB}(\boldsymbol{\nu}) + \mathbf{C}_A(\boldsymbol{\nu}) \in \mathbb{R}^{6\times 6}, \\
\mathbf{D}(\boldsymbol{\nu}) &= \mathbf{D}_L + \mathbf{D}_{NL}(\boldsymbol{\nu}) \in \mathbb{R}^{6\times 6}, \\
\mathbf{g}(\boldsymbol{\eta}) &\in \mathbb{R}^{6}.
\end{aligned}
\end{beqn}
The rigid-body terms depend only on the vehicle itself. $\mathbf{M}_{RB}$ is the
rigid-body mass and inertia matrix, and $\mathbf{C}_{RB}(\boldsymbol{\nu})$
contains the rigid-body Coriolis and centripetal terms. The hydrodynamic
contributions come from the surrounding water. The added-mass matrix
$\mathbf{M}_A$ accounts for the inertia of the fluid that is accelerated
together with the vehicle, and $\mathbf{C}_A(\boldsymbol{\nu})$ contains the
corresponding added-mass Coriolis and centripetal terms. The damping matrix
$\mathbf{D}(\boldsymbol{\nu})$ combines linear skin friction $\mathbf{D}_L$ with
quadratic drag $\mathbf{D}_{NL}(\boldsymbol{\nu})$, which grows with the
velocity. Finally, $\mathbf{g}(\boldsymbol{\eta})$ is the vector of restoring
forces and moments from gravity and buoyancy, which act at the centre of
gravity and the centre of buoyancy respectively. Placing the centre of gravity
below the centre of buoyancy gives the vehicle passive stability in roll and
pitch.

The control forces and moments are produced by the thrusters,
\begin{beqn}
\begin{aligned}
\boldsymbol{\tau} = \mathbf{T}\mathbf{f}, \qquad
\mathbf{T} \in \mathbb{R}^{6\times 8}, \quad \mathbf{f} \in \mathbb{R}^{8},
\end{aligned}
\end{beqn}
where $\mathbf{f}$ holds the forces of the eight thrusters on the BlueROV2 Heavy
and $\mathbf{T}$ is the thruster configuration matrix, determined by the
position and direction of each thruster. Since
$\operatorname{rank}(\mathbf{T}) = 6$, the vehicle is fully actuated in all six
degrees of freedom. With eight thrusters for six degrees of freedom, it is in
fact over-actuated, so the mapping from $\boldsymbol{\tau}$ to $\mathbf{f}$
has redundancy.
